
\documentclass[12pt, reqno]{amsart}
\usepackage{amsmath, amsthm, amscd, amsfonts, amssymb, graphicx, xcolor}
\usepackage{tikz}
%\usetikzlibrary{arrows.meta, decorations.markings}
\usepackage[bookmarksnumbered, colorlinks, plainpages]{hyperref}
\usepackage{tikz}
\usepackage{enumitem}
\usepackage{bm}
\usepackage{float}
%
\textheight 22.5truecm \textwidth 14.5truecm
\setlength{\oddsidemargin}{0.35in}\setlength{\evensidemargin}{0.35in}
%
\setlength{\topmargin}{-.5cm}

\newtheorem{theorem}{Theorem}[section]
\newtheorem{lemma}[theorem]{Lemma}
\newtheorem{proposition}[theorem]{Proposition}
\newtheorem{corollary}[theorem]{Corollary}
\theoremstyle{definition}
\newtheorem{definition}[theorem]{Definition}
\newtheorem{example}[theorem]{Example}
\newtheorem{exercise}[theorem]{Exercise}
\newtheorem{conclusion}[theorem]{Conclusion}
\newtheorem{conjecture}[theorem]{Conjecture}
\newtheorem{criterion}[theorem]{Criterion}
\newtheorem{summary}[theorem]{Summary}
\newtheorem{axiom}[theorem]{Axiom}
\newtheorem{observation}[theorem]{Observation}
\newtheorem{problem}[theorem]{Problem}
\theoremstyle{remark}
\newtheorem{remark}[theorem]{Remark}
\numberwithin{equation}{section}
\newtheorem{notation}{Notation}[section]
\DeclareMathOperator{\Hom}{Hom}
\DeclareMathOperator{\im}{im}



\begin{document}

\setcounter{page}{1}

%
%\color{darkgray}{
%\noindent 
%{\small Annals of Mathematics and Computer Science}\hfill     {\small ISSN: 2789-7206}\\
%{\small Vol 20 (2024) 1-3}\hfill  {\small https://doi.org/10.56947/amcs.v20.223}}


\centerline{}

\centerline{}


%------------------------------------------------------------------------------

%Title of the paper
\title[]{Maximum number of  cells in a matroid subdivision of hypersimplex }

%Author names and affiliations
\author[]{ Mahyar Hasani}



\footnote{This work was partially supported by INSF under Grant No. 4049675.}
%\address{$^{1}$ Department of Mathematics,University of Münster, Münster , GERMANY.}
%\email{\textcolor[rgb]{0.00,0.00,0.84}{lutz.hille@uni-muenster.de}}




%\address{$^{2}$ Department of Mathematics, IPM - Institute for Research in Fundamental Sciences, Tehran, Iran}
%\email{\textcolor[rgb]{0.00,0.00,0.84}{hadi@ipm.ir}}



\address{$^{3}$ Department of Mathematics, IPM - Institute for Research in Fundamental Sciences, Tehran, Iran}
\email{\textcolor[rgb]{0.00,0.00,0.84}{mahyarhasani@ipm.ir}}


%\dedicatory{This paper is dedicated to Professor ABCD}
%
%\date{Received: xxxxxx; Revised: yyyyyy; Accepted: zzzzzz.
%\newline \indent $^{*}$ Corresponding author
%\newline \indent © The Author(s) 2025. This article is licensed under a Creative Commons Attribution-
%\newline \indent NonCommercial-NoDerivatives 4.0
%International License. To view a copy of the licence, visit 
%\newline \indent \url{https://creativecommons.org/licenses/by-nc-nd/4.0/}}
%

%Abstract, keywords, math subject classification
\begin{abstract}

David Speyer proved that
the maximum number of  cells in a  regular matroid subdivision of the hypersimplex
$\Delta(m,t)$
is
$\binom{t-2}{m-1}$
, using the language of tropical geometry and properties of the Tutte polynomial of a matroid \cite[Theorem 6.1]{Speyer}.
In this paper, we prove this result in general, i.e. for arbitrary matroid subdivisions including the non-regular ones. We perform  a different  combinatorial way using our results in \cite{HHH} and  independent of tropical language. 
\newline
\newline
\noindent \textit{Keywords.} Algebraic Geometry, Matroid Theory, Hypersimplex, Inverse Limit, Tropical Geometry.    
\newline
%\noindent \textit{2020 Mathematics Subject Classification.} Primary 46L55; Secondary 44B20
\end{abstract} \maketitle



%------------------------------------------------------------------------------



\tableofcontents



\section{Introduction}







Triangulations and subdivisions are among  the most important notions in convex geometry as well as  in algebraic geometry and topology. For example in toric geometry,  we can associate an algebraic  variety to a triangulation of a special polytope.   More generally we can consider face-fitting decompositions, in which we do not need the 
smaller polytopes to be simplices anymore. One of the most important special cases is a subdivision of hypersimplex 

\begin{center}
	$\Delta(m,t)= 
	\left\{(x_{1},x_{2},...,x_{t})\in \mathbb{R}^t:x_{1}+...+x_{t}=m, 0\le x_{i}\le 1\right\}$  
\end{center}

by matroid polytopes. In fact, the hypersimplex is itself the polytope of the uniform matroid $U_{m,t}$ and its subdivision is also considered in matroid theory\cite{oriented}.
On the other hand one can construct toric varieties from 
matroid polytopes
involved in a matroid subdivision and  attach them to get the so called stable varieties or "shas"(stable hyperplane arrangements) as explained in \cite{Alex}. A natural question is what is the maximum number of irreducible components in a "shas" obtained by a matroid subdivision of hypersimplex 
$\Delta(m,t)$?
Or, in another words, what 
is the maximum number of polytopes in a matroid subdivision of the hypersimplex? We will recall some definitions from matroid theory:

\begin{definition}
	A \emph{matroid} is a pair 
	$P=(\left\{1,...,t\right\},\varGamma)$,
	where 
	$\varGamma$
	is  a collection  
	of subsets of
	$\left\{1,...,t\right\}$
	satisfying the properties below:
	
	
	
	\begin{enumerate}
		\item 
		$ \emptyset \in \varGamma $
		\item 
		If $I\in \varGamma$ and $J\subseteq I$, then $J\in \varGamma$
		\item 
		If $I,J \in \varGamma$  and $|I| < |J|$ , there exists $j\in J-I$
		such that $I\cup  \left\{j\right\} \in  \varGamma $.
		
		
	\end{enumerate}
	
	We also denote the matroid simply by $\Gamma$.
	
\end{definition}



\begin{definition}
	If the  pair 
	$P=(\left\{1,...,t\right\},\varGamma)$ is a matroid, 
	the elements of $\varGamma$ are  called independent sets of the matroid.
	
\end{definition}


Matroids can also be defined via their bases:

\begin{definition}
	Let the  pair 
	$P=(\left\{1,...,t\right\},\varGamma)$ be a matroid, a subset 
	$I\subseteq \left\{1,...,t\right\}$ is called a \emph{basis} if it is a  
	maximal independent set. We denote the set of all bases by $\beta(P)$ or $\beta(\Gamma)$.
	
\end{definition}




%\begin{definition} 
%	Assume that t vectors
%	$f_1,...,f_t$
%	span 
%	$\mathbb{C}^m$ as a $\mathbb{C}$-vector space. A subset $I\subseteq \left\{1,...,t\right\}$
%	is called \emph{independent} with respect to these vectors if the set $\left\{f_i:i\in I\right\}$	is linearly independent over $\mathbb{C}$.
%\end{definition}

%\begin{definition}
%	Assume again that vectors 
%	$f_1,...,f_t$
%	span 
%	$\mathbb{C}^m$. The \emph{linear matroid} represented by the vectors 
%	$f_1,...,f_t$ in $\mathbb{C}^m$ is the pair 
%	$P=(\left\{1,...,t\right\},\varGamma)$,
%	where 
%	$\varGamma$
%	is the set of all independent sets. The number $m$ is called the \emph{rank} of $P$.
%\end{definition}


%\begin{definition}
%	If 
%	$P=(\left\{1,...,t\right\},\varGamma)$
%	is a linear matroid represented by vectors 
%	$f_1,...,f_t$ in $\mathbb{C}^m$, a subset 
%	$I\subseteq \left\{1,...,t\right\}$ is called a \emph{basis} if the vectors 
%	$\left\{f_i:i\in I\right\}$
%	form a basis for $\mathbb{C}^m$. We denote the set of all bases by $\beta(P)$.
%\end{definition}



\begin{definition}
	The \emph{matroid polytope} of a matroid $P=(\left\{1,...,t\right\},\varGamma)$ is the convex hull of the points $e_I$ with $I\in \beta(P)$. It is a polytope in $\mathbb{R}^t$ and is denoted by $M_P$ or $M_\Gamma$. 
\end{definition}




Two operations can be performed on matroids, one is \emph{deletion} and the other is \emph{contraction} :

\begin{definition}\label{del}
	Let $P=(\left\{1,...,t\right\},\varGamma)$ be a matroid and $X\subseteq \left\{1,...,t\right\}$. Put 
	\begin{center}
		$\varGamma|X=\left\{ I\subseteq X: I\in \varGamma\right\}$
	\end{center}
	
	Then the pair 
	$(X,\varGamma|X)$ is also a matroid, called the \emph{deletion of}$X^c$  and is denoted by $P|X$.
	
\end{definition}


\begin{definition}\label{cont}
	Let $P=(\left\{1,...,t\right\},\varGamma)$ be a matroid and $X\subseteq \left\{1,...,t\right\}$.
	
	Put 
	
	\[
	\varGamma/X = \left\{ I \subseteq X^c: P|X \text{ has a basis } B \text{ such that } B \cup I \in \varGamma \right\}
	\]
	
	Then the pair 
	$(X^c,\varGamma/X)$ is also a matroid, called the \emph{contraction} of $X$  and is denoted by $P/X$.
	
	
	%I(M/T) = {I ⊆E−T :I ∪BT ∈I(M)}
\end{definition}



The notion of  rank function for a matroid is essential in our study:

\begin{definition}
	Let $P=(\left\{1,...,t\right\},\varGamma)$ be a matroid and $X\subseteq \left\{1,...,t\right\}$. The rank of $X$, denoted by $r_{P}(X)$, is defined to be the cardinality of a basis of $P|X$.
\end{definition}



We will work with two kinds of matroids:

\begin{definition}
		A  matroid $P=(\left\{1,...,t\right\},\varGamma)$ of rank $m$ is \emph{linear} if there are
		  vectors 
		 	$f_1,...,f_t$ in $\mathbb{C}^m$
		 	such that  a subset
		 	$I\subseteq \left\{1,...,t\right\}$
		 	is independent if and only if
		 	the set $\left\{f_i:i\in I\right\}$	is linearly independent over $\mathbb{C}$.
\end{definition}
 There is a notion of connectedness in matroid theory which is useful for our aim.

\begin{definition}
	A  matroid $P=(\left\{1,...,t\right\},\varGamma)$  is
	connected if one cannot decompose the ground set 
\begin{center}
		$\left\{1,...,t\right\}=E_1\cup E_2$
\end{center}
	such that
	 a subset
	$I\subseteq \left\{1,...,t\right\}$
	is independent if and only if
	
	\begin{center}
		$I=I_1\cup I_2$
	\end{center}
	where
	$I_1\subseteq E_1$
	,
		$I_2\subseteq E_2$
	 and $I_1$, $I_2$
	are independent.
	
	
\end{definition}


%\begin{definition}
%	If $P=(\left\{1,...,t\right\},\varGamma)$ is a linear matroid represented by vectors 
%	$f_1,...,f_t$ in $\mathbb{C}^m$ and $I\subseteq \left\{1,...,t\right\}$, the \emph{rank} of $I$, denoted by $r_P(I)$ is the dimension of the subspace generated by  vectors $f_i$ for $i\in I$.
%\end{definition}



\begin{definition}
	A collection of maximal matroid polytopes $\left\{Q_j: j\in J\right\}$ in $\mathbb{R}^t$ is a \emph{matroid decomposition} of $\Delta(m,t)$
	if $\bigcup\limits_{j\in J} Q_j=\Delta(m,t)$ and they intersect face-fittingly; i.e., for each  pair  $Q_i$ and $Q_j$ , $Q_i\cap Q_j$ is either empty or a face of both $Q_i$ and $Q_j$.
\end{definition}




It is proven by 
Herrmann,  Jensen, Joswig, and Sturmfels
that 
the maximum number of  cells in a matroid subdivision of the hypersimplex
$\Delta(3,t)$ is   not bounded by a constant \cite{Tropical}.  Valery Alexeev also proved that 
this number is bounded by $(t-3)^2$  for $\Delta(3,t)$ \cite{Alex}. 
Finally, David Speyer proved that for all $m,t$   
the maximum number of  cells in a regular matroid subdivision of 
$\Delta(m,t)$
is
$\binom{t-2}{m-1}$.
 In this article we extend this
result to all arbitrary(including non-regular)  matroid subdivisions of 
$\Delta(m,t)$ via the ideas written in 
\cite{HHH}. We hope our new approach also helps us to answer a question about  the Dressian fan  in tropical geometry, denoted by $Dr(m,t)$. We briefly state the question below:

 Every element in the Dressian fan
 $Dr(m,t)$ determines a matroid subdivision for
 $\Delta(m,t)$. The number, \emph{spread} of an element in the Dressian, is the 
 number of maximal cells in the corresponding matroid subdivision of hypersimplex. Question 35 in \cite{Dressian} states:
 \begin{itemize}
 	\item 
 What is the maximal spread of a ray of $Dr(m,t)$ (at least for m = 3)?
   \end{itemize}
  
  
  In the question above, rays  correspond to  subdivisions that
  cannot be coarsened, i.e. 
  if you want to coarsen it, you will only get the whole hypersimplex $\Delta(m,t)$.
  We hope that the recursive process that we used for the previous problem can also be used to solve the problem above.
  
  
\section{Relation to inverse limit of point configurations}

The notion of inverse limit for  
an action of an algebraic group on an algebraic variety was studied in \cite{Th99} using Geometric Invariant Theory \cite{Mumford}
and its variation \cite{Th94,DH98}. For a reference concerning this matter  and the special case of point configurations see
\cite{BH}.



\begin{definition}
	A point configuration with parameters $m,t$ is a tuple $(p_1,...,p_t)$ such that 
	$p_i\in \mathbb{P}^{m-1}$ for every $i=1,...,t$. We also assume that we have at least $m+1$ general points in this configuration.
\end{definition}


One can consider finitely many quotients of the natural action of $PGL(m)$ on $(\mathbb{P}^{m-1})^t$, provided by 
Geometric Invariant Theory, and consider their inverse limit.  In fact, for every weight 
$\theta \in \Delta(m,t)$, we have a corresponding moduli space 
$M^{\theta}$ parametrizing $\theta$-semistable representations up to S-equivalence. Based on the Geometric Invariant Theory literature, they form an inverse system and we denote this inverse limit by 
$\overline{M}(m,t)$.

 In fact, the walls 
$W^l_I$ for 
$I\subseteq \left\{1,...,t\right\}$ , $1\leq l<m$  decompose the hypersimplex 
$\Delta(m,t)$ into finitely many chambers. For each chamber in this decomposition, we have a distinct moduli space.
We can specify a representative $\theta$ for each chamber; this representative determines the corresponding moduli space $M^\theta$.
The inverse limit
$\overline{M}(m,t)$
is a closed subvariety of
$\prod\limits_{\theta\in C(m,t)}  M^\theta$,
where 
$C(m,t)$  is the set which has exactly one weight as a representative from each chamber. In this way, we can denote each element of the inverse limit 
$\overline{M}(m,t)$
as 
$(\varphi_{\theta})_{\theta\in C(m,t)}$ such that $\varphi_{\theta}\in M^{\theta}$.


We need to establish some notations:

\begin{notation}
	
	
	Let $e_1,...,e_t$ be the standard basis vectors of $\mathbb{R}^t$. Suppose $I\subseteq \left\{1,...,t\right\}$, $0\leq l<m$ and $\theta=(\theta_1,...,\theta_t) \in \mathbb{R}^t$.
	\begin{center}
		\begin{enumerate}
			\item 
			$e_I=\sum\limits_{i\in I} e_{i}$ 
			
			
			\item 
			$\theta(I)=\theta_I=\sum\limits_{i\in I} \theta_{i}$ 
			\item 
			$W^l_I=\left\{(\theta_1,...,\theta_t)\in \mathbb{R}^t : \sum\limits_{i\in I} \theta_{i}=l \right\}$ 
		%	\item 
		%	$W_I^{\le l}=\left\{(\theta_1,...,\theta_t)\in \mathbb{R}^t : \sum\limits_{i\in I} \theta_{i}\leq l \right\}$
			
		%	\item 
		%	$W_I^{\ge l}=\left\{(\theta_1,...,\theta_t)\in \mathbb{R}^t : \sum\limits_{i\in I} \theta_{i}\geq l \right\}$
			
			
		%	\item  For a vector $\vec{d}=(d_0,...,d_{t})$, we put 
		%	$\theta(d)=\sum\limits_{i=1}^t\theta_i d_i$.
		\end{enumerate}
	\end{center}
\end{notation}


In the above notations, the  particular case  $I=\left\{i \right\}$ is of interest to us. In this case  we have

\begin{center}
	\begin{itemize}
		\item 
			$W^0_{\left\{i\right\}} \cong \Delta(m,t-1)$
			\item 
		$W^1_{\left\{i\right\}} \cong \Delta(m-1,t-1)$
		\end{itemize}
\end{center}
We will work with both of them multiple times.


  In the paper 
\cite{HHH}, we construct algebraic morphisms 


\begin{center}
	$
	\overline{M}(m, t) \longrightarrow \overline{M}(m, t-1)
	$
	
	
	$
	\overline{M}(m, t) \longrightarrow \overline{M}(m-1, t-1)
	$
	
	
\end{center}

for every 
$1 \leq i\leq t$
which are called
 \emph{deletion} and \emph{contraction} maps. They are
 analogous 
 to the well-known
 maps 
 of the
  Kapranov-Chow quotients of Grassmannians \cite{Kap1}. One can combine the 
   \emph{deletion} and \emph{contraction} maps
   for a fixed 
   $1 \leq i\leq t$
    simultaneously
    to get the natural morphism
   \begin{center}
   	$
   	\overline{M}(m, t) \longrightarrow \overline{M}(m, t-1)\times \overline{M}(m-1, t-1)
   	$.
   	\end{center}
   
  
  
 If we have a point of the
  inverse limit
  	$
  \overline{M}(m, t)$, this gives us a matroid decomposition of the hypersimplex
  $\Delta(m,t)$  for $m=2,3$ by \cite[Theorem 1.14]{HHH}, we can apply the morphisms above to 
  get  elements in 
  	$
  \overline{M}(m, t-1)$
  and 
  	$
  \overline{M}(m-1, t-1)$ and decompositions for
  $\Delta(m,t-1)$ and $\Delta(m-1,t-1)$ respectively. 
  
  In fact these morphisms
  are derived from natural maps:
  
  \begin{center}
  	
  	$\eta_i:\prod\limits_{\theta\in C(m,t)}  M^\theta\longrightarrow \prod\limits_{\theta\in C(m,t-1)}  M^\theta$
  	
  	$
  	(\varphi_{\theta})_{} \mapsto (\varphi_{\theta}^{D_i})_{}
  	$
  			
  	\end{center}
  	
  	\begin{center}
  
  	$\gamma_i:\prod\limits_{\theta\in C(m,t)}  M^\theta\longrightarrow \prod\limits_{\theta\in C(m-1,t-1)}  M^\theta$
  	
  	$
  	(\varphi_{\theta})_{} \mapsto (\varphi_{\theta}^{C_i})_{}
  	$
  \end{center}
  where
  $\varphi_{\theta}^{D_i}$ and
  $\varphi_{\theta}^{C_i}$
  are obtained from the representation 
  $\varphi_{\theta}$ by applying the 
  \emph{deletion} and \emph{contraction} process over point $i$
   for each
  $1 \leq i\leq t$  as explained in 
  \cite[Definition
  3.14,Definition 3.15]{HHH}. 
   In this definition, 
  chambers of the hypersimplices
  $\Delta(m,t-1)$ and $\Delta(m-1,t-1)$
  are obtained with intersecting of chambers of 
  $\Delta(m,t)$ with 
  $W^0_{\left\{i\right\}}$,$W^1_{\left\{i\right\}}$ respectively. 
  So if we have a point of the inverse limit   
  $(\varphi_{\theta})_{\theta \in C(m,t)}$
  in
  $
  \overline{M}(m, t)$
   the images would be tuples of the form
   $(\varphi_{\theta}^{D_i})_{\theta\in C(m,t-1)}$
   and
   $(\varphi_{\theta}^{C_i})_{\theta\in C(m-1,t-1)}$ respectively.
  
  
  
  For 
  a  point inside the  inverse limit
  $\overline{M}(m,t)$,
   we denote the
  maximum number of cells in the corresponding  subdivision of hypersimplex
  $\Delta(m,t)$ by K(m,t). Then we can get a recursive formula if we prove the items below:
 % \begin{lemma}\label{keyitems}
  
 
  \begin{enumerate}
  	\item 
  	Every component of 
  	 $(\varphi_{\theta}^{D_i})_{\theta\in C(m,t-1)}$
  	and
  	$(\varphi_{\theta}^{C_i})_{\theta\in C(m-1,t-1)}$ is obtained from one component of 
  	 $(\varphi_{\theta})_{\theta \in C(m,t)}$.
  	

  	\item 
  		Every  component of the 
  		initial element of
  		$(\varphi_{\theta})_{\theta \in C(m,t)}$
  		in $\overline{M}(m,t)$ appears in at least one(or both) of the components of the images
  		$(\varphi_{\theta}^{D_i})_{\theta\in C(m,t-1)}$
  		and
  		$(\varphi_{\theta}^{C_i})_{\theta\in C(m-1,t-1)}$
  		in
  	$\overline{M}(m,t-1)$ and 
  	$\overline{M}(m-1,t-1)$ respectively.
  \end{enumerate}
  % \end{lemma}

  If we prove these two facts, then we have:
  
  \begin{center}
  	
  	\begin{itemize}
  	%	\item 
  	%	$K(m,t-1)\leq K(m,t)$
  	%%	$K(m-1,t-1)\leq K(m,t)$
  		\item 
  		$K(m,t-1)+K(m-1,t-1)\geq K(m,t)$
  		
  	\end{itemize}
  	  \end{center}
  	
  	To see why this is true, take a  tuple with 
  	$K(m,t)$ maximal cells
  	 which is  a point of the inverse limit
  	$\overline{M}(m,t)$, the items (1), (2) in the previous paragraph show that 
  	we can decompose the cells of this tuple into union of two sets $A\cup B$ such that $A$  and $B$ are the
  	maximal cells of the images in 
  	$\overline{M}(m,t-1)$
  	and
  	$\overline{M}(m-1,t-1)$
  	respectively, so
  	
  	\begin{center}
  		
  		$K(m,t)=|A\cup B|\leq |A|+|B|$
  	\end{center}
  	
  	
  	
  	 Thus using induction on $m+t$ and Pascal identity, we get
  	 
  	 \begin{center}
  	 	
  	 $K(m,t)\leq \binom{t-3}{m-1}+ \binom{t-3}{m-2}=\binom{t-2}{m-1}$
  	
  	\end{center}

   for 
   $m\geq t+1$ and $m\geq 2$ and the bound is achieved. The base of induction is also clear ($m=2,t=4$).
   
 % \begin{observation}
  %The case (1) above is clear by definition. For the case (2),
  %let 
  %$(P_s)^q_{s=1}$
  %is a tuple which is an element of the inverse limit 
  %$\overline{M}(m,t)$. 
  % Consider a fixed  point configuration
  %$P_s=(p_1,...,p_t)$ in this tuple where $p_i\in \mathbb{P}^{m-1}$. For each point $p_i$ we have two cases, either there is another point $p_j$ such that the position of $p_i$ and $p_j$ has  coincidence or for every $j\neq i$
%  the position of $p_j$ is different from $p_i$.\begin{itemize}
 % 	\item 
 %If $p_i$ and $p_j$ coincide then $M_{P_s}$ has a non empty intersection with  the region
 %$W^0_{\left\{j\right\}}$,
  % because if we select 
 %$\theta\in 	M_{P_s}$
  % we can  define
   %$\theta^\prime=(\theta^\prime_1,...,\theta^\prime_t)$
   %as below:
   %
   %
   %\[ 
   %\theta^\prime_k = 
   %\begin{cases}
   	
   	%\theta_{i}+\theta_{j}& k=j  \\
   	%0 & k=i \\
   	%\theta_{k}	& k\neq  i,j
   %	
   %\end{cases}
   %\]
   
   
   
   
   
   %In this way, $\theta^\prime \in M_{P_s}\cap W^0_{\left\{j\right\}}$ and if you remove the point $p_i$ from the %point configuration $P_s$ then $P_s^{D_i}$ would appear in the image
    %$\overline{M}(m,t-1)$
    %as a maximal component. 
  
  
  
  %\item 
  %If the position of $p_i$ is different from every
  %$p_j$ for $j\neq i$.
  
  
  % In the same way as the previous one, we can put 
  %\[ 
  %\theta^\prime_k = 
  %\begin{cases}
  %	
  %	1& k=i  \\
  %	\theta_{k}-\frac{1-\theta_{i}}{t-1}	& k\neq  i
  %\end{cases}
 % \]
  %One can see 
  %$ M_{P_s}\cap W^1_{\left\{j\right\}}\neq \emptyset$, because in this case, we are increasing the  weight of $i$ %-th point and decreasing
   %weights of other points. Thus there are some weights in 
   %$W^1_{\left\{j\right\}}$
    %such that still satisfies the necessarily inequalities of $M_{P_s}$. More precisely, choose a basis 
   %$\left\{f_s: s\in \beta\right\}$  for $\mathbb{R}^m$ where
   %$\beta\subseteq\left\{1,2,...,t\right\}$
   %and  has  $f_i$ as one of vectors and also has $m-1$ vectors choosen from 
   %outside of 
 %$  \left\{i\right\}$.  Then 
 %it is clear that $e_\beta\in M_{P_s}$ and 
 %$e_\beta( \left\{i\right\})=e_i=1$
  % thus 
  % $e_\beta\in M_{P_s}\cap W^1_{\left\{i\right\}}$.
   %So
   %if you project over the point $p_i$, in the point configuration $P_s$ then $P_s^{C_i}$ would appear in the image
   %$\overline{M}(m-1,t-1)$
  %as a maximal component.
   
   % more precisely the only inequality that may be violated is that the $i-$th point might lie in a  $l-$dimensional sub-linear space with  points $\left\{p_j:j\in J\right\}$   
  %where
  %$2\leq l< m$
  %and
  %$J\subseteq \left\{1,2,...,t\right\}-\left\{i\right\}$ and $\theta^\prime$ should satisfies the corresponding relation. Since $\theta \in M_{P_i}$
  %we have:
  
  %\begin{center}
  %	
  %	$
  %	\theta_i+\theta_J\leq l
  %	$
  %\end{center}
  
  %The following calculations show that this inequality still holds for $\theta^\prime$.
  
  
  %\begin{center}
  	
  %	$
  %	\theta^\prime_i+\theta^\prime_J= 1+\sum\limits_{j\in J} \theta_j^\prime=1+\sum\limits_{j\in J} \theta_j - \frac{|J|}{t-1}(1-\theta_i)=1+\theta_J- \frac{|J|}{t-1}(1-\theta_i)\leq \theta_i+\theta_J- \frac{|J|}{t-1}(1-\theta_i) \leq l
  %	$
  %\end{center}
  
  
  
  
  
  
  
  
%\end{itemize}
  %\end{observation}
  
  
  It is quite easy to see that both of (1), (2) are satisfied for point configurations or linear matroids. 
  % but now we give just a motivation that why it should be true for linear matroids:
  With this motivation we prove (1),(2) for general matroids in theorem \ref{main} and state this bound for every matroid decomposition of hypersimplex
  $\Delta(m,t)$
  independent of inverse limit  and regularity concept. 
  
  
 %	What 
  %David Speyer
  %showed is valid only for regular matroid subdivisions,
  %we want to show this   for arbitrary matroid subdivision of hypersimplex $\Delta(m,t)$. 
  
  
  
  \section{Proof that the bound is working }
  We first  establish some notation and lemmas, specially the lemma below, which is well-known in matroid theory.
  \begin{lemma}\label{connecteddim}
  	Let
  	$(\left\{1,...,t\right\},\varGamma)$
  	be an abstract matroid,
  	the matroid polytope
  	$M_{\Gamma}$
  	is of maximal dimension if and only if the matroid $\Gamma$ is connected.
  \end{lemma}
  
  
  \begin{proof}
  	Dimension of the matroid polytope can be computed using \cite[Proposition 2.4]{DMP},i.e
  	\begin{center}
  		$\dim(M_{\Gamma})=t-C(\Gamma)$
  	\end{center}
  	where
  	$C(\Gamma)$
  	is the number of connected components
  	of the matroid
  	$\Gamma$. Notice that a matroid  $\Gamma$ is connected if and only  of $C(\Gamma)=1$. Since the matroid polytope $M_{\Gamma}$ is a sub-polytope of the hypersimplex $\Delta(m,t)$ and $\dim(\Delta(m,t))=t-1$ the result is driven.
  \end{proof}
  
  
  \begin{notation}
  		If 
  		$P=(\left\{1,...,t\right\},\varGamma)$ is a matroid and $1\leq i\leq t$, we can operate the 
  	 \emph{Deletion} and \emph{Contraction} process , over point $i$ and get two new matroids
  	  $\Gamma^{D_i}$ and $\Gamma^{C_i}$
  	  with indices
  	  $\left\{1,2,..., i-1,i+1,...t\right\}$. 
  	 
  \end{notation}
  
  %\begin{lemma}
  %	If $P=(p_1,...,p_t)$ is a point configuration and $1\leq i\leq t$,
  %	then
  %	
  %	\begin{itemize}
  %%		$M_{P}\cap (\theta_i=0)=M_{P^{D_i}}$
  	%	\item 
  	%	$M_{P}\cap (\theta_i=1)=M_{P^{C_i}}$
  	%	
  	%\end{itemize}
  %
  	%We want to show that for every matroid subdivision of hypersimplex $\Delta(m,t)$
  %	
  %\end{lemma}
  
  We also need  the definition of isomorphism between matroids:
  
  \begin{definition}\label{isomat}
  	Two abstract matroids 
  	$(\left\{1,...,t\right\},\varGamma)$ and 
  	$(\left\{1,...,t\right\},\varGamma^\prime)$
  	are isomorphic if there is a bijective map 
  	$f:\left\{1,...,t\right\}\to \left\{1,...,t\right\}$
  	such that 
  	for every subset
  	$I\subseteq\left\{1,...,t\right\}$, 
  	$I\in \varGamma$
  	if and only if 
  	$f(I)\in \varGamma^\prime$.
  \end{definition}
  
  \begin{definition}
  	
  	
  	We say two matroid subdivisions
  	
  	\begin{center}
  	$\left\{M_{\Gamma_1},...,M_{\Gamma_q}\right\}$	
  		
  	\end{center}
  
  
  \begin{center}
  	,
  \end{center}	
  	
  	
  	
  	
  	
  	\begin{center}
  		$\left\{M_{\Gamma^\prime_1},...,M_{\Gamma^\prime_q}\right\}$	
  		
  	\end{center}
  	
  	
  	
  	
  	
  	
  	
  	 are isomorphic if there is a bijection between their cells, i.e
  	 there are isomorphisms of matroids
  	 	$f_i:\Gamma_i\to \Gamma_i^\prime$
  	  for
  	  $1\leq i\leq q$.
  	
  	
  \end{definition}
  
  Let $D(m,t)$ be the set of all matroid subdivisions of the hypersimplex $\Delta(m,t)$ up to the notion of isomorphism defined in the above. We will construct the map  below 
  analogous the inverse limit one:
  
  
  \begin{center}
  	\[
  D(m, t)  \longrightarrow  D(m, t-1)\times D(m-1, t-1)
  	\]
  	
  	\[
  	(M_{\Gamma_s})^q_{s=1}
  	\longmapsto
  (	(M_{\Gamma_s}\cap \Delta(m,t-1))^q_{s=1}, 	(M_{\Gamma_s}\cap \Delta(m-1,t-1))^q_{s=1})
  	\]
  \end{center}
  
  
  Here we should notice that since the decomposition 
  	$(M_{\Gamma_s})^q_{s=1}\in D(m,t)$
  is face-fitting, after intersecting with
  $\Delta(m,t-1)$ and $\Delta(m-1,t-1)$, it will remain face-fitting.
  We should also prove that both of the components 
  $M_{\Gamma_s}\cap \Delta(m,t-1)$ and $M_{\Gamma_s}\cap \Delta(m-1,t-1)$ are matroid polytopes. We will prove this in the proposition below:
  
  
  %\begin{notation}
  %	If
  %	$(\left\{1,...,t\right\},\varGamma)$
  %is an abstract matroid and $1\leq i\leq t$, we can operate %the 
  %	\emph{Deletion} and \emph{Contraction} process for this matroid, over point $i$ and get two new matroids
  %	$P^{D_i}$ and $P^{C_i}$
  %	with indices
  %	$\left\{1,2,..., i-1,i+1,...t\right\}$
  %	and 
  %	$\left\{j:\pi_i(p_j)=0\right\}$ respectively.  
  %	
  %\end{notation}
  
  
 
  
  
  
  
  \begin{proposition}\label{intersection}
  	
  	If 
  	$(\left\{1,...,t\right\},\varGamma)$
  	 is an abstract matroid  such that $M_{\Gamma}$ is of maximal dimension and $1\leq i\leq t$,
  	then
  	
  	\begin{itemize}
  		\item 
  		$M_{\Gamma}\cap W^0_{\left\{i\right\}}\cong M_{\Gamma^{D_i}}$
  		\item 
  		$M_{\Gamma}\cap W^1_{\left\{i\right\}}\cong M_{\Gamma^{C_i}}$
  		
  	\end{itemize}
  	
  \end{proposition}

\begin{proof}
	Defenition of matroid polytope 
	$M_{\Gamma}$
	says:
	\begin{center}
			$M_{\Gamma}\cap W^0_{\left\{i\right\}}=conv\left\{e_I:I\in \beta(\Gamma), i\notin I \right\}$
			
				$M_{\Gamma}\cap W^1_{\left\{i\right\}}=conv\left\{e_I:I\in \beta(\Gamma), i\in I \right\}.$
			
	\end{center}
	Based on the assumption the matroid polytope $M_{\Gamma}$ is of maximal dimension. Thus lemma \ref{connecteddim} implies that 
	the matroid
	$\Gamma$ is connected and the element $i$ is neither loop nor coloop. So we have
		\begin{center}
		$\beta(\Gamma^{D_i})=\left\{I\in \beta(\Gamma): i\notin I \right\}$
		
		$\beta(\Gamma^{C_i})=\left\{I-\left\{i\right\}:I\in \beta(\Gamma), i\in I \right\}$
		
	\end{center}
	 Now it is easy to see the isomporphisms.
\end{proof}
 Now if we prove
items (1), (2) of the theorem below , then we get the required bound by the same way as before.
 \begin{theorem}\label{main}
	With the notations above, we have:
	
	\begin{enumerate}
		\item 
		Every  maximal component of
		$(M_{\Gamma_s}\cap \Delta(m,t-1))^q_{s=1}$
		and
		$(M_{\Gamma_s}\cap \Delta(m-1,t-1))^q_{s=1}$
		 is obtained from one component of 
		$(M_{\Gamma_s})^q_{s=1}$.
		
		
		\item 
		Every   component of the 
		initial decomposition 
			$(M_{\Gamma_s})^q_{s=1}$
		in $D(m,t)$ appears in at least one of the maximal components of the images
			$(M_{\Gamma_s}\cap \Delta(m,t-1))^q_{s=1}$
		and
		$(M_{\Gamma_s}\cap \Delta(m-1,t-1))^q_{s=1}$
		in
		$D(m,t-1)$ and 
		$D(m-1,t-1)$ respectively.
	\end{enumerate}
\end{theorem}

\begin{proof}
	
	Item (1) is clear. For (2) notice that   the intersection 
	$M_{\Gamma_s}\cap \Delta(m,t-1)$
	and
	$M_{\Gamma_s}\cap \Delta(m-1,t-1))$
	could  be non-maximal, in this case they are not  computed as a maximal component. We should prove  that for every component 
	$M_{\Gamma_s}$ at least one of 
	$M_{\Gamma_s}\cap \Delta(m,t-1)$
	or
	$M_{\Gamma_s}\cap \Delta(m-1,t-1)$
	is maximal. This is proven in the next proposition.
	
	\end{proof}
  
  
  
  
  
  \begin{proposition}
  	At least one of the polytopes
  	 $M_{\Gamma_s}\cap \Delta(m,t-1)$
  	or
  	$M_{\Gamma_s}\cap \Delta(m-1,t-1)$
  	is maximal dimensional.
  \end{proposition}
  
  \begin{proof}
  	Based of proposition \ref{intersection} we have to prove at least one of the polytopes
  	$M_{\Gamma_s^{D_i}}$ or 
  	$M_{\Gamma_s^{C_i}}$
  	is of maximal dimensional. By lemma \ref{connecteddim} we should prove that 
  	either
  		$\Gamma_s^{D_i}$ or 
  	$\Gamma_s^{C_i}$
  	 is connected. Since
  	 	$(M_{\Gamma_s})^q_{s=1}\in D(m,t)$,
  	  the matroid polytope $M_{\Gamma_s}$ is of maximal dimension, so 
  	$\Gamma_s$ is connected. As known in matroid theory, Theorem 4.3.1 in Oxley book \cite{Ox} implies that either 
  	  $\Gamma^{D_i}$ or 
  	  $\Gamma^{C_i}$
  	  is connected.
  \end{proof}
  
  
  
  \begin{corollary}
  Maximum number of maximal cells in an arbitrary matroid decomposition of hypersimplex $\Delta$(m,t) is
  $\binom{t-2}{m-1}$.
  	
  \end{corollary}
  

%
%Again, note how we refer to Theorem \ref{theo1} and formula \eqref{2.1}.
%\\
%\\

{\bf Acknowledgement.} 

I would like to express my gratitude to my advisors, Prof. Lutz Hille and Prof. S. Mohammad Hadi Hedayatzadeh, for their constant support and invaluable guidance. In particular, I am grateful to Prof. Lutz Hille for introducing me to this interesting problem and the concept of inverse limit, which played a key role in shaping my research.


\bibliographystyle{amsplain}
\begin{thebibliography}{99}

\bibitem{Speyer} D. Speyer, Tropical linear spaces, SIAM Journal on Discrete Mathematics 22 (2008), no. 4, 1527–1558. \url{https://arxiv.org/pdf/2005.01787}



\bibitem{BH} M. Blume and L. Hille. Quivers and moduli spaces of pointed curves of genus zero, Algebr. Comb., 4(1):89–124,
 2021
\url{https://doi.org/10.5802/alco.152}



%\bibitem{TevSch}
 %L. Schaffler and J. Tevelev, Compactifications of Moduli of Points and Lines in the Projective Plane,
%International Mathematics Research Notices 2022 (2021), no. 21, 17000–17078.
%\url{https://doi.org/10.1093/imrn/rnab200}

\bibitem{DH98} I. Dolgachev, Y. Hu, Variation of geometric invariant theory quotients, Publ. Math.
IHES 87 (1998), 5–51
\url{https://doi.org/10.1007/BF02698859}





\bibitem{Th94} M. Thaddeus, Geometric invariant theory and flips. Journal of the American Mathematical Society, vol. 9, no. 3, pp. 691-723, 1996. 
\url{https://arxiv.org/pdf/alg-geom/9405004}



\bibitem{Th99} M. Thaddeus, Complete collineations revisited. Math. Ann. 315 (1999), 469–495,
\url{https://doi.org/10.1007/s002080050324}




\bibitem{Kap1} M. Kapranov, Chow quotients of Grassmannians I, Adv. Soviet Math. 16 (1993),
29–110,
\url{https://arxiv.org/abs/alg-geom/9210002}


%\bibitem{GP}Lothar Gerritzen and Meinolf Piwek. Degeneration of point configurations in the projective plane.
%Indag. Math. (N.S.) 2 (1991), no. 1, 39–56,
%\url{https://www.sciencedirect.com/science/article/pii/001935779190040E}


%\bibitem{kap2}  M.M. Kapranov, B. Sturmfels, A.V Zelevinsky: Quotients of toric varieties. Math. Ann. 290
%(1991), no. 4, 643–655.
%\url{https://doi.org/10.1007/BF01459264}


%\bibitem{Hu} Yi Hu. Topological aspects of Chow quotients. Journal of Differential Geometry,
%69(3):399–440, 2005.\url{https://arxiv.org/pdf/math/0308027}



%\bibitem{Baker} Hendrik Baker, Jurgen Hausen, and Simon Keicher. On Chow quotients of torus actions.
%Michigan Math. J., 64(3):451–473, 2015.
%\url{https://arxiv.org/pdf/1203.3759}




 %\bibitem{matpol} I.M.Gel'fand, V.V.Serganova, Combinatorial geometries and torus strata on homogeneous compact manifolds,  Russian Mathematical Surveys, 1987. 
%\url{https://people.bath.ac.uk/masadk/zzz/gelf-serg87.pdf}





 \bibitem{Alex} Valery Alexeev, Moduli of Weighted Hyperplane Arrangements, Advanced Courses in Mathematics. CRM Barcelona, Springer
, Basel, 2015.
\url{https://doi.org/10.1007/978-3-0348-0915-3}






\bibitem{Ox} James Oxley. Matroid theory, Oxford
University Press, 2011.
\url{https://www.math.ens.psl.eu/~benoist/refs/Oxley.pdf}.



\bibitem{Mumford} D. Mumford, J. Fogarty, F. Kirwan: Geometric invariant theory. Third edition. Ergebnisse
der Mathematik und ihrer Grenzgebiete (2), 34. Springer-Verlag, Berlin, 1994.
\url{https://link.springer.com/book/9783540569633}

%\bibitem{HKT} Paul Hacking, Sean Keel, and Jenia Tevelev, Compactification of the moduli space of hyperplane arrangements,
%Journal of Algebraic Geometry 15 (2006), no. 4, 657–680.
%\url{https://doi.org/10.1090/S1056-3911-06-00445-0}


%\bibitem{King} A.D. King, Moduli of representations of finite dimensional algebras, Quart. J. Math. Oxford (2) 45(1994) 515–530.
%\url{https://www.math.uni-bielefeld.de/~sek/sem/stability/king.pdf}

 %\bibitem{KeeTeveq} Sean Keel, Jenia Tevelev. Equations for M0,n. Int. J. Math., 20(09):1159–1184,
 %2009.
 
 %\url{https://www.semanticscholar.org/paper/Equations-for-%24%5Cbar-M_%7B0%2Cn%7D%24-Keel-Tevelev/f530d40827a3f578e815e0c2fb77cfe0abdd78b6}
 
 
 %\bibitem{GHP} L. Gerritzen, F. Herrlich, M. van der Put, Stable n-pointed trees of projective
 %lines, Nederl. Akad. Wetensch., Indag. Math. 50 (1988), 131–163.
 %\url{https://pure.rug.nl/ws/files/3373983/1988MathProcAGerritzen.pdf}

%\bibitem{gruber1989facet}P. Gruber, S.Ry{\v{s}}kov. %Facet-to-facet implies face-to-face. European Journal of Combinatorics. 10
%(1989), 83–84. 
%\url{https://www.sciencedirect.com/science/article/pii/S0195669889800356}.




 %\bibitem{Liu} Qing Liu. Algebraic Geometry and Arithmetic Curves. Oxford University Press, pp.576, 2006.
%\url{https://www.gbv.de/dms/goettingen/337835292.pdf}


%\bibitem{Hoskins} V.Hoskins, Moduli problems and geometric invariant theory, Lecture notes 2016. vol.7. no.28. p. 12,2015.





%\bibitem{Dolgachev} V.Dolgachev, Lectures on invariant theory, 2003.   no.296, Cambridge University.





\bibitem{Tropical}Herrmann, Jensen, Joswig, Sturmfels,
How to draw tropical planes, Electron. J. Combin. 16 (2009), n

o. 2,
Special volume in honor of Anders Bjorner, Research Paper 6, 26.
\url{https://arxiv.org/pdf/0808.2383}

\bibitem{Dressian} S.Herrmann, M.Joswig, D.Speyer, Dressians, tropical Grassmannians, and their rays. Forum Math. 26(6), 2014.
\url{https://arxiv.org/pdf/1112.1278}

\bibitem{oriented} M. Celaya, G.Loho, C. Yuen, Oriented matroids from triangulations of products of simplices, Selecta Mathematical, 2024, Springer.
\url{https://arxiv.org/pdf/2005.01787}



\bibitem{HHH} M. Hedayatzadeh, M. Hasani, Inverse limit and Chow quotient of point configurations in projective
spaces are not isomorphic, preprint, 2026.

\bibitem{DMP} E. Feichtner, B. Sturmfels, Matroid polytopes, nested sets and Bergman
fans, Port. Math. (N.S.), 62 (2005), pp. 437–468.





%
%\bibitem{Bekka1} Bekka, B. (2006). Property (T) for C$^*$-algebras. Bulletin of the London Mathematical Society, 38(5), 857-867.
%\url{https://doi.org/10.1112/S0024609306018765}
%
%\bibitem{Bekka2}  Bekka, B., de La Harpe, P., \& Valette, A. (2008). Kazhdan's property (T). Cambridge university press.
%\url{https://doi.org/10.1017/CBO9780511542749}
%
%\bibitem{Brown} Brown, N. P. (2006). Kazhdan's property T and C$^*$-algebras. Journal of Functional Analysis, 240(1), 290-296.
%\url{https://doi.org/10.1016/j.jfa.2006.05.003}

%
%\bibitem{Connes}  Connes, A., \& Jones, V. (1985). Property T for von Neumann algebras. Bulletin of the London Mathematical Society, 17(1), 57-62.
%\url{https://doi.org/10.1112/blms/17.1.57}
%
%\bibitem{Green}  Green, P. (1978). The local structure of twisted covariance algebras. Acta Mathematica, 140, 191-250.
%\url{https://doi.org/10.1007/BF02392308}
%    
%\bibitem{Kazhdan} Kazhdan, D. A. (1967). Connection of the dual space of a group with the structure of its close subgroups. Funktsional'nyi Analiz i ego Prilozheniya, 1(1), 71-74.
%\url{https://doi.org/10.1007/BF01075866}
%
%\bibitem{Leung} Leung, C. W., \& Ng, C. K. (2009). Property (T) and strong property (T) for unital C$^*$-algebras. Journal of Functional Analysis, 256(9), 3055-3070.
%\url{https://doi.org/10.1016/j.jfa.2009.01.004}

\end{thebibliography}

\end{document}

%------------------------------------------------------------------------------
% End of journal.tex
%------------------------------------------------------------------------------
